\section{Results: Completed Development Evidence}
The scheduled main study has completed all three sources and all three
training seeds: 558 evaluation cells, independently audited from
protocol-v2 listing outcomes. This is retrospective development evidence,
not a new confirmatory test. Follow-up tuning and direct-comparator
audits have separate completion status and are reported within their
verified scopes below.
Earlier development-v1 numbers are withdrawn as deployability evidence
because that implementation used a private simulation mask for query
eligibility. Protocol v2 corrects that information error.

\begin{table*}[t]
\centering
\caption{Complete three-source development at 10\% capacity. Inflation is
the mean positive log-quote gap, averaged over three training seeds;
parentheses give its sample standard deviation across seeds. Price errors
use each source's native price units and are not pooled across sources.
Methods within each source share the evaluation cohort.}
\label{tab:synthetic_development}
\begin{tabular}{llrrr}
\toprule
Source & Method & Inflation (SD) & Honest MAE & Attack MAE\\
\midrule
Synthetic & Mask + benefit & .04963 (.00186) & 3503.62 & 4242.44\\
 & Robust, no feedback & .01339 (.00119) & 3568.22 & 3634.80\\
 & SAV + benefit output & .01928 (.00333) & 3546.75 & 3752.06\\
 & SAV-Val & .02075 (.00167) & 3522.46 & 3754.03\\
\midrule
JUCars & Mask + benefit & .01933 (.00243) & 6875.38 & 6883.06\\
 & Robust, no feedback & .01154 (.00044) & 7081.70 & 7078.81\\
 & SAV + benefit output & .01811 (.00209) & 7059.94 & 7046.86\\
 & SAV-Val & .01549 (.00128) & 7084.77 & 7073.94\\
\midrule
MUCars & Mask + benefit & .04612 (.00089) & 27281.98 & 28282.98\\
 & Robust, no feedback & .00774 (.00108) & 27578.04 & 27793.87\\
 & SAV + benefit output & .00753 (.00044) & 26958.21 & 27124.56\\
 & SAV-Val & .00689 (.00030) & 27147.41 & 27311.95\\
\bottomrule
\end{tabular}
\end{table*}

On the synthetic market, full SAV-Val has 58.19\% less positive inflation than the mask-augmented
predictor with ordinary-benefit acquisition. Its honest MAE is 0.54\%
higher. That comparison alone could suggest progress, but the matched
and stronger controls give a different assessment. SAV-Val has 7.63\%
more inflation than its ordinary-benefit-output variant. The paired
reference-minus-full difference is $-0.00147$, with a conditional
95\% batch interval $[-0.00394,0.00104]$. Thus the primary component
contrast does not establish an advantage for the suppression output in
this synthetic setting.

The robust/no-feedback control has the lowest inflation among these
four synthetic methods. SAV-Val's inflation is 54.94\% higher, while its honest
MAE is 1.28\% lower. The paired reference-minus-full inflation difference
is $-0.00736$, with interval $[-0.00982,-0.00514]$. Attacked MAE also
favors the simpler control. These results support a limited trade-off,
not an overall improvement by adding the full feedback mechanism.
The interval uses the ten shared evaluation batches after averaging
paired contrasts across fitted seeds; it is not a confidence interval
over new training datasets or real seller populations.

Remaining unilateral regret is zero for both full SAV-Val and the
robust/no-feedback control in the saved synthetic 10\% cells, within
the numerical tolerance used by the independent audit. This concerns
the returned finite-game reports and unilateral moves only. It is not
a proof of a unique or globally optimal equilibrium. Other controls
have nonzero remaining regret and nonconverged batches, so their
reported outcomes should also be considered alongside the saved
unilateral envelopes rather than interpreted as exact worst-case
solutions.

The JUCars primary component contrast favors the suppression-output
variant: SAV-Val has 14.45\% less inflation than its ordinary-benefit
counterpart, with honest MAE 0.35\% higher. The reference-minus-full
difference is $0.00262$, with conditional batch interval
$[0.00140,0.00379]$. This source-specific result does not reverse the
synthetic finding or establish an advantage over the simpler robust
control. Relative to robust/no-feedback, SAV-Val has 34.23\% more
inflation; the reference-minus-full difference is $-0.00395$, with
interval $[-0.00558,-0.00249]$. Its honest MAE is 0.04\% higher.
The attacked-MAE means differ only slightly (Table~\ref{tab:synthetic_development});
the trade-off must remain visible rather than attributing all differences
to the action score.

MUCars gives a different comparison. Full SAV-Val has 10.92\% less
inflation than robust/no-feedback and 1.56\% lower honest MAE. The
reference-minus-full inflation difference is $0.000845$, with conditional
batch interval $[0.000383,0.001242]$. Its ordinary-benefit component
contrast is also positive: 8.45\% less inflation, with difference
$0.000636$ and interval $[0.000426,0.000896]$, while honest MAE is
0.70\% higher. This component improvement is below the prespecified
10\% magnitude criterion. The complete main study therefore supports
source-dependent gains, not a uniform advantage of feedback or of the
suppression output. At the primary capacity, all three MUCars SAV-Val
and robust/no-feedback fits have zero remaining unilateral regret within
the audit tolerance; this remains a finite-game diagnostic.

\subsection{Original robust predictor with ordinary-benefit acquisition}
The saved robust/no-feedback control in the main study uses the suppression
output. To complete the comparison specified in the original contract, we
separately evaluate the same nine frozen robust predictors and their
multi-output, out-of-fold deployment heads using the ordinary-benefit
output. All six capacities, source cohorts, seeds and seller-search settings
are retained. The 54 additional cells are independently checked for public
scores, allocations, replay, metrics and exact first-batch deviations.
Honest quotes are identical to the suppression-output robust control;
this follow-up changes acquisition and the induced seller reports, without
refitting or selecting models from evaluation outcomes.

\begin{table*}[t]
\centering
\caption{Completed robust/ordinary-benefit reference at 10\% capacity.
Reference values average three training seeds, with sample standard deviation
in parentheses. The last column is reference minus full SAV-Val inflation,
with a 95\% conditional shared-batch interval; positive values favor SAV-Val.
Price errors retain native source units.}
\label{tab:robust_benefit_reference}
\begin{tabular}{lrrrr}
\toprule
Source & Reference inflation (SD) & Honest MAE & Attack MAE & Difference [95\% interval]\\
\midrule
Synthetic & 0.01077 (0.00066) & 3568.22 & 3613.74 & $-0.00998\;[-0.01260,-0.00791]$\\
JUCars & 0.01117 (0.00046) & 7081.70 & 7081.51 & $-0.00432\;[-0.00598,-0.00285]$\\
MUCars & 0.00778 (0.00108) & 27578.04 & 27770.09 & $0.00089\;[0.00045,0.00130]$\\
\bottomrule
\end{tabular}
\end{table*}

At the primary capacity, SAV-Val has 92.69\% more inflation than this
reference on the synthetic cohort and 38.65\% more on JUCars. Their
reference-minus-full intervals are below zero. On MUCars, SAV-Val has
11.42\% less inflation, with difference $0.000888$ and interval
$[0.000445,0.001295]$; honest MAE is 1.56\% lower and attacked MAE
is 1.65\% lower. This meets the original fixed-setting comparison's
10\% reduction, positive interval and at most 5\% honest-error deterioration
criteria on MUCars only. It does not meet them across all scheduled sources,
and the different component criterion in protocol v2 remains separate.
The completed equal-allowance extension separately tunes and refits both
static acquisition references and their deployment heads across all nine
source/size cohorts. These original 54 fixed-setting cells and intervals remain
unchanged; the separate 324-cell extension is reported below. Neither study
establishes a universal feedback advantage.


\subsection{All preset verification capacities}
Figure~\ref{fig:capacity} retains all six capacities fixed before this
run. The comparison with robust/no-feedback depends on both source
and capacity. On the synthetic cohort, SAV-Val has 11.92\% less
inflation at 30\% capacity; the reference-minus-full difference is
$0.00136$, with an exploratory conditional batch interval
$[0.00085,0.00202]$. At 20\%, the difference is close to zero and its
interval crosses zero. On JUCars, full SAV-Val has higher inflation
than the robust control at every preset capacity, including 30\%.
On MUCars, the reference-minus-full inflation difference is positive
at all six capacities, including zero, with conditional intervals above
zero. The zero-query difference is consistent with an estimator-level
effect and cannot be attributed to acquiring additional fields.
These are unadjusted descriptive frontier comparisons. They do not
replace the prespecified 10\% primary comparison or establish a
general transition point.

\begin{figure*}[!t]
\centering
\includegraphics[width=\textwidth]{figures/capacity_frontiers.pdf}
\caption{All preset capacity curves for the completed three-seed cohorts.
Top: source means and sample standard deviations across three training
seeds. Bottom: robust/no-feedback minus SAV-Val inflation, with unadjusted
95\% shared-batch intervals conditional on the fitted models; positive
differences favor SAV-Val. Paired contrasts are averaged over seeds before
resampling the ten shared batches. The dotted vertical line marks the
prespecified 10\% primary capacity. These retrospective development
frontiers include all three scheduled sources.}
\label{fig:capacity}
\end{figure*}

\subsection{Restricted numerical SMI domain diagnostic}
Separate from the full-field study, we fit deterministic linear SVRs
on year, mileage, and fiscal power (synthetic/MUCars) or seats (JUCars). Only mileage
can be withheld; there are no verification queries. Full and reduced
models use identical estimator-training rows with recorded mileage,
training-only feature and log-target normalization, $C=1$, and
$\epsilon=0.15$. The two-stage SMI equations~\cite{zhang2020strategic}
are evaluated with free, normalized nonnegative, and normalized unit-box
missing-coordinate domains. Strict consistency is required for all
retained coefficient signs. This convention and the domains are declared
adaptations, not an official Weka reproduction. Known validation inputs
are not clipped to the training range. We enumerate both seller reports
exactly and independently check the LP outputs with an analytic
one-coordinate solution. The full/reduced fits use 3,356 synthetic,
2,677 JUCars, and 3,500 MUCars records respectively.

\begin{table*}[t]
\centering
\caption{Restricted one-field numerical diagnosis: one deterministic fit
per source and 1,000 evaluation records. MAE is in native source price
units. Outside is the fraction of imputed mileage coordinates outside the
estimator-training normalized range. There were no solver fallbacks in
these completed cells. These methods are not directly compared with the
full-field, shared-budget SAV-Val study.}
\label{tab:numeric_smi}
\begin{tabular}{llrrr}
\toprule
Source & Method/domain & Honest MAE & Inflation & Outside\\
\midrule
Synthetic & Median SVR & 9328.99 & .02868 & 0\\
 & SMI / free & 9621.74 & 0 & 1\\
 & SMI / nonnegative & 9621.74 & 0 & 1\\
 & SMI / unit box & 9398.70 & $2.18\!\times\!10^{-6}$ & 0\\
\midrule
JUCars & Median SVR & 8568.44 & .01911 & 0\\
 & SMI / free & 9045.97 & 0 & 1\\
 & SMI / nonnegative & 8595.09 & 0 & 0\\
 & SMI / unit box & 8595.09 & 0 & 0\\
\midrule
MUCars & Median SVR & 49767.58 & .00394 & 0\\
 & SMI / free & 49767.58 & 0 & 1\\
 & SMI / nonnegative & 49767.58 & 0 & 1\\
 & SMI / unit box & 49767.58 & 0 & 0\\
\bottomrule
\end{tabular}
\end{table*}

Table~\ref{tab:numeric_smi} shows why variable domains and prediction
error must accompany incentive metrics. Free-domain SMI eliminates
positive inflation on these cohorts but increases honest MAE on
synthetic and Jordanian records. On JUCars, its mean signed log bias is
$-0.44095$, versus $-0.20506$ for the nonnegative domain and $-0.18022$
for median SVR. The free and nonnegative synthetic imputations all lie
outside the training range. Unit-box SMI has one profitable synthetic
withholding report among 1,000 listings, with mean inflation
$2.18\times10^{-6}$. Its recorded normalized mileage is $1.00403$;
withholding permits an imputation of $1$ under a negative mileage
coefficient. A bounded imputation domain is not a guarantee
for an unbounded evaluation-input domain. These restricted observations
do not establish a general SMI guarantee or a full-market SAV-Val
advantage.

\subsection{Fixed-estimator counterfactual mixture development}
The fixed-estimator CDM study has completed all nine source/seed pairs.
The initial seven-pair generation and its two originally unavailable MUCars
pairs retain separate immutable provenance; their complete evidence union
has an independent audit of 1,512 execution cells and 108 calibration fits.
Every stored payoff row, full-matrix LP optimum, rule allocation and mixture
expectation is checked. Table~\ref{tab:cdm_development} includes all three
scheduled sources and all three fitting seeds.

\begin{table*}[t]
\centering
\caption{CDM development at 10\% capacity with fixed estimators. Entries are
three-seed means of positive expected-log uplift. The final column is the
conditional shared-batch 95\% interval for risk minus CDM; positive values
favor CDM. This compares verification rules within each fixed estimator,
rather than a matched full-estimator comparison with feedback-trained SAV.}
\label{tab:cdm_development}
\begin{tabular}{llrrrrl}
\toprule
Source & Estimator & CDM & Risk & Benefit & Uniform & Risk $-$ CDM interval\\
\midrule
Synthetic & Mask-aware & 0.06224 & 0.05634 & 0.06152 & 0.07968 & $[-0.00779,-0.00385]$\\
 & Robust & 0.01192 & 0.01005 & 0.01017 & 0.01340 & $[-0.00291,-0.00110]$\\
\midrule
JUCars & Mask-aware & 0.01879 & 0.01854 & 0.01953 & 0.01886 & $[-0.00042,-0.00012]$\\
 & Robust & 0.01063 & 0.01055 & 0.01126 & 0.01009 & $[-0.00044,0.00029]$\\
\midrule
MUCars & Mask-aware & 0.04398 & 0.03102 & 0.04605 & 0.04430 & $[-0.01461,-0.01144]$\\
 & Robust & 0.00699 & 0.00626 & 0.00782 & 0.00658 & $[-0.00102,-0.00045]$\\
\bottomrule
\end{tabular}
\end{table*}

On the synthetic cohort, CDM reduces uplift relative to uniform mixing but
has higher uplift than the simple risk rule under both estimators. On JUCars,
it improves on benefit priority but does not establish a risk-rule advantage;
under the robust estimator it also loses to uniform mixing. All twelve pure
rules, the training-selected pure rule, both mixtures and all six capacities
are retained in the accompanying contrast data. These comparisons with uniform mixing do not establish an advantage
over the stronger risk rule.

The complete MUCars replication also favors risk priority over CDM under
both price predictors: CDM uplift is 41.79\% and 11.63\% higher for
mask-aware and robust prices. The corresponding risk-minus-CDM intervals
are below zero. With robust prices, CDM improves on benefit priority by
10.62\% and on the training-selected pure rule by 6.55\%, but has 6.32\%
more uplift than uniform mixing. At the preset 30\% capacity, robust-price
MUCars instead has 22.19\% less CDM uplift than risk priority, with paired
difference $0.00136$ and exploratory interval $[0.00107,0.00169]$.
This capacity-dependent result does not replace the primary comparison.
A finite calibration improvement does not imply a general held-out average
improvement over simple alternatives.

\begin{figure*}[t]
\centering
\includegraphics[width=\textwidth]{figures/cdm_capacity_frontiers_full.pdf}
\caption{Complete CDM source study at every preset capacity. Top:
mask-aware price predictor; bottom: robust/no-feedback price predictor.
Points average three fitting seeds and bars give their sample standard
deviations. The 10\% primary capacity is marked. All sources and capacities
remain visible; no favorable rule or source is selected after evaluation.
Conditional risk-minus-CDM primary intervals appear in
Table~\ref{tab:cdm_development}.}
\label{fig:cdm_capacity}
\end{figure*}

The finite LP guarantee concerns the stored calibration types and peer
profiles. It does not guarantee a reduction of mean held-out uplift, and its
worst-gain objective differs from that average. In the source experiment,
the honest-error constraint is invariant across rules and contributes no
observed discrimination. The mechanism's usefulness therefore remains an
empirical question despite exact finite optimization.

\subsection{Equal-allowance tuning and size retraining}
The independently audited follow-up retains 9 complete three-seed source/cap cohorts and 972 final evaluation cells. All nine scheduled source/cap cohorts and all 972 evaluation cells are independently audited. Table~\ref{tab:equal_tuning_development} retains the three tuned methods at the primary capacity. Fixed-weight versions and every preset capacity remain in the associated audit and contrast outputs.

\begin{table*}[t]
\centering
\caption{Every completed three-seed source/size cohort after equal three-trial allowances,
at 10\% verification capacity. Means use three training seeds and the same
1,000 evaluation listings at every cap. Tuning selects honest plus attacked
squared log loss inside training folds, rather than minimizing uplift alone.
All nine scheduled source/cap cohorts are complete and independently audited.}
\label{tab:equal_tuning_development}
\begin{tabular}{lrlrrr}
\toprule
Source & Training cap & Tuned method & Inflation & Honest MAE & Attack MAE\\
\midrule
Synthetic & 1000 & Robust, no feedback & 0.01121 & 3937.99 & 4033.31\\
 &  & SAV + benefit output & 0.01469 & 3942.56 & 4120.77\\
 &  & SAV-Val & 0.01921 & 3902.46 & 4137.45\\
\midrule
Synthetic & 2500 & Robust, no feedback & 0.01386 & 3606.63 & 3669.80\\
 &  & SAV + benefit output & 0.01348 & 3601.19 & 3706.48\\
 &  & SAV-Val & 0.01790 & 3600.74 & 3753.64\\
\midrule
Synthetic & 5000 & Robust, no feedback & 0.01910 & 3530.14 & 3685.92\\
 &  & SAV + benefit output & 0.01433 & 3551.61 & 3659.43\\
 &  & SAV-Val & 0.01809 & 3541.78 & 3745.14\\
\midrule
JUCars & 1000 & Robust, no feedback & 0.00680 & 6962.61 & 6958.75\\
 &  & SAV + benefit output & 0.00550 & 6961.08 & 6956.38\\
 &  & SAV-Val & 0.00538 & 6959.43 & 6954.92\\
\midrule
JUCars & 2500 & Robust, no feedback & 0.01079 & 6985.52 & 6976.39\\
 &  & SAV + benefit output & 0.00865 & 6977.20 & 6970.41\\
 &  & SAV-Val & 0.00866 & 6966.07 & 6961.63\\
\midrule
JUCars & 5000 & Robust, no feedback & 0.00826 & 7008.73 & 7000.31\\
 &  & SAV + benefit output & 0.01244 & 7041.13 & 7030.07\\
 &  & SAV-Val & 0.01097 & 7003.56 & 6994.50\\
\midrule
MUCars & 1000 & Robust, no feedback & 0.00901 & 31573.54 & 31623.98\\
 &  & SAV + benefit output & 0.00846 & 31478.84 & 31562.14\\
 &  & SAV-Val & 0.00830 & 31644.91 & 31682.27\\
\midrule
MUCars & 2500 & Robust, no feedback & 0.00865 & 28126.15 & 28425.48\\
 &  & SAV + benefit output & 0.00841 & 27828.02 & 27978.40\\
 &  & SAV-Val & 0.00698 & 27663.99 & 27806.81\\
\midrule
MUCars & 5000 & Robust, no feedback & 0.00715 & 27444.54 & 27656.87\\
 &  & SAV + benefit output & 0.00633 & 26860.26 & 26939.22\\
 &  & SAV-Val & 0.00564 & 27021.60 & 27150.16\\
\bottomrule
\end{tabular}
\end{table*}

The complete Synthetic cohort uses configured cap 1,000. Against the tuned static suppression-output control, inflation is 71.34\% higher, with reference-minus-SAV interval $[-0.00993,-0.00593]$. Honest and attacked MAE changes are -0.90\% and +2.58\%; their paired error intervals exclude zero. Against the tuned feedback benefit-output variant, inflation is 30.75\% higher, with reference-minus-SAV interval $[-0.00648,-0.00244]$. Honest and attacked MAE changes are -1.02\% and +0.40\%; their paired error intervals include zero in at least one metric.

The complete Synthetic cohort uses configured cap 2,500. Against the tuned static suppression-output control, inflation is 29.20\% higher, with reference-minus-SAV interval $[-0.00609,-0.00209]$. Honest and attacked MAE changes are -0.16\% and +2.28\%; their paired error intervals include zero in at least one metric. Against the tuned feedback benefit-output variant, inflation is 32.82\% higher, with reference-minus-SAV interval $[-0.00556,-0.00338]$. Honest and attacked MAE changes are -0.01\% and +1.27\%; their paired error intervals include zero in at least one metric.

The complete Synthetic cohort uses configured cap 5,000. Against the tuned static suppression-output control, inflation is 5.26\% lower, with reference-minus-SAV interval $[-0.00111,0.00310]$. Honest and attacked MAE changes are +0.33\% and +1.61\%; their paired error intervals include zero in at least one metric. Against the tuned feedback benefit-output variant, inflation is 26.21\% higher, with reference-minus-SAV interval $[-0.00658,-0.00109]$. Honest and attacked MAE changes are -0.28\% and +2.34\%; their paired error intervals include zero in at least one metric.

The complete JUCars cohort uses configured cap 1,000. Against the tuned static suppression-output control, inflation is 20.94\% lower, with reference-minus-SAV interval $[0.00100,0.00186]$. Honest and attacked MAE changes are -0.05\% and -0.05\%; their paired error intervals include zero in at least one metric. Against the tuned feedback benefit-output variant, inflation is 2.19\% lower, with reference-minus-SAV interval $[-0.00027,0.00047]$. Honest and attacked MAE changes are -0.02\% and -0.02\%; their paired error intervals include zero in at least one metric.

The complete JUCars cohort uses configured cap 2,500. Against the tuned static suppression-output control, inflation is 19.79\% lower, with reference-minus-SAV interval $[0.00157,0.00280]$. Honest and attacked MAE changes are -0.28\% and -0.21\%; their paired error intervals include zero in at least one metric. Against the tuned feedback benefit-output variant, inflation is 0.16\% higher, with reference-minus-SAV interval $[-0.00082,0.00078]$. Honest and attacked MAE changes are -0.16\% and -0.13\%; their paired error intervals include zero in at least one metric.

The complete JUCars cohort uses configured cap 5,000. Against the tuned static suppression-output control, inflation is 32.71\% higher, with reference-minus-SAV interval $[-0.00368,-0.00157]$. Honest and attacked MAE changes are -0.07\% and -0.08\%; their paired error intervals include zero in at least one metric. Against the tuned feedback benefit-output variant, inflation is 11.83\% lower, with reference-minus-SAV interval $[0.00095,0.00197]$. Honest and attacked MAE changes are -0.53\% and -0.51\%; their paired error intervals exclude zero.

The complete MUCars cohort uses configured cap 1,000. Against the tuned static suppression-output control, inflation is 7.92\% lower, with reference-minus-SAV interval $[0.00050,0.00093]$. Honest and attacked MAE changes are +0.23\% and +0.18\%; their paired error intervals include zero in at least one metric. Against the tuned feedback benefit-output variant, inflation is 1.91\% lower, with reference-minus-SAV interval $[-0.00004,0.00036]$. Honest and attacked MAE changes are +0.53\% and +0.38\%; their paired error intervals exclude zero.

The complete MUCars cohort uses configured cap 2,500. Against the tuned static suppression-output control, inflation is 19.25\% lower, with reference-minus-SAV interval $[0.00116,0.00225]$. Honest and attacked MAE changes are -1.64\% and -2.18\%; their paired error intervals exclude zero. Against the tuned feedback benefit-output variant, inflation is 16.98\% lower, with reference-minus-SAV interval $[0.00116,0.00171]$. Honest and attacked MAE changes are -0.59\% and -0.61\%; their paired error intervals exclude zero.

The complete MUCars cohort uses configured cap 5,000. Against the tuned static suppression-output control, inflation is 21.00\% lower, with reference-minus-SAV interval $[0.00101,0.00191]$. Honest and attacked MAE changes are -1.54\% and -1.83\%; their paired error intervals exclude zero. Against the tuned feedback benefit-output variant, inflation is 10.76\% lower, with reference-minus-SAV interval $[0.00033,0.00103]$. Honest and attacked MAE changes are +0.60\% and +0.78\%; their paired error intervals include zero in at least one metric.

Equal allowances do not establish a consistent full-method advantage. Tuning minimizes honest plus attacked squared log loss; a chosen setting can improve accuracy while increasing inflation. No objective, trial grid, primary capacity or reference identity was changed after observing outcomes. These conditional, unadjusted intervals concern fitted models and historical validation batches. Every scheduled source/cap cohort, negative finding and honest-error tradeoff is retained.

\subsection{Ordinary-benefit reference with equal selection allowances}
The separate extension evaluates ordinary-benefit acquisition with a static
robust estimator at the same three gap-weight trials and twelve residual
epochs. Static candidate price fits can be shared with the suppression-output
control, but the ordinary-benefit validation objective is recomputed and its
selected weight is refitted when necessary. Its outer predictions and final
action head follow its own selections. The comparison therefore does not
relabel a suppression-output policy or select a reference on evaluation data.

\begin{table*}[t]
\centering
\caption{Every fully completed three-seed source/cap cohort in the audited
ordinary-benefit allowance extension, at 10\% capacity. BEN is the static robust
predictor with ordinary-benefit acquisition. Gain is $100(\Phi_{BEN}-\Phi_{SAV})/
\Phi_{BEN}$. Error changes are SAV relative to BEN; positive values mean larger
MAE. Intervals concern the paired reference-minus-SAV inflation difference,
averaged over seeds before resampling ten shared batches. Fixed means
$\beta=0.5$; tuned selects from the same three training-internal trials.}
\label{tab:equal_benefit_reference}
\small
\begin{tabular}{lrlrrrrrr}
\toprule
Source & Cap & Role & SAV $\Phi$ & BEN $\Phi$ & Gain (\%) & Difference interval & Honest $\Delta$ (\%) & Attack $\Delta$ (\%)\\
\midrule
JUCars & 1000 & Fixed & 0.00605 & 0.00596 & -1.51 & $[-0.00050,0.00026]$ & -0.23 & -0.26\\
 &  & Tuned & 0.00538 & 0.00675 & +20.31 & $[0.00087,0.00189]$ & -0.05 & -0.05\\
\midrule
JUCars & 2500 & Fixed & 0.00846 & 0.00993 & +14.72 & $[0.00072,0.00218]$ & -0.84 & -0.86\\
 &  & Tuned & 0.00866 & 0.01125 & +23.02 & $[0.00200,0.00308]$ & -0.28 & -0.24\\
\midrule
JUCars & 5000 & Fixed & 0.01549 & 0.01117 & -38.65 & $[-0.00593,-0.00286]$ & +0.04 & -0.11\\
 &  & Tuned & 0.01097 & 0.00863 & -27.09 & $[-0.00333,-0.00125]$ & -0.07 & -0.09\\
\midrule
MUCars & 1000 & Fixed & 0.00847 & 0.00802 & -5.56 & $[-0.00062,-0.00028]$ & -0.01 & -0.10\\
 &  & Tuned & 0.00830 & 0.00879 & +5.68 & $[0.00032,0.00068]$ & +0.23 & +0.20\\
\midrule
MUCars & 2500 & Fixed & 0.00857 & 0.00880 & +2.61 & $[-0.00021,0.00067]$ & -1.59 & -1.79\\
 &  & Tuned & 0.00698 & 0.00861 & +18.92 & $[0.00112,0.00222]$ & -1.64 & -2.10\\
\midrule
MUCars & 5000 & Fixed & 0.00689 & 0.00778 & +11.42 & $[0.00045,0.00128]$ & -1.56 & -1.65\\
 &  & Tuned & 0.00564 & 0.00709 & +20.41 & $[0.00097,0.00186]$ & -1.54 & -1.78\\
\midrule
Synthetic & 1000 & Fixed & 0.02095 & 0.00936 & -123.82 & $[-0.01280,-0.01038]$ & -1.58 & +2.92\\
 &  & Tuned & 0.01921 & 0.01258 & -52.62 & $[-0.00868,-0.00453]$ & -0.90 & +1.44\\
\midrule
Synthetic & 2500 & Fixed & 0.01764 & 0.00926 & -90.60 & $[-0.01029,-0.00662]$ & -0.90 & +2.46\\
 &  & Tuned & 0.01790 & 0.01220 & -46.79 & $[-0.00774,-0.00368]$ & -0.16 & +2.33\\
\midrule
Synthetic & 5000 & Fixed & 0.02075 & 0.01077 & -92.69 & $[-0.01268,-0.00786]$ & -1.28 & +3.88\\
 &  & Tuned & 0.01809 & 0.01391 & -30.09 & $[-0.00682,-0.00204]$ & +0.33 & +3.14\\
\bottomrule
\end{tabular}
\end{table*}

At JUCars cap 1,000, tuned SAV-Val has 20.31\% lower inflation than the tuned ordinary-benefit robust reference; the reference-minus-SAV interval favors SAV-Val.
At JUCars cap 2,500, tuned SAV-Val has 23.02\% lower inflation than the tuned ordinary-benefit robust reference; the reference-minus-SAV interval favors SAV-Val.
At JUCars cap 5,000, tuned SAV-Val has 27.09\% higher inflation than the tuned ordinary-benefit robust reference; the reference-minus-SAV interval favors the reference.
At MUCars cap 1,000, tuned SAV-Val has 5.68\% lower inflation than the tuned ordinary-benefit robust reference; the reference-minus-SAV interval favors SAV-Val.
At MUCars cap 2,500, tuned SAV-Val has 18.92\% lower inflation than the tuned ordinary-benefit robust reference; the reference-minus-SAV interval favors SAV-Val.
At MUCars cap 5,000, tuned SAV-Val has 20.41\% lower inflation than the tuned ordinary-benefit robust reference; the reference-minus-SAV interval favors SAV-Val.
At Synthetic cap 1,000, tuned SAV-Val has 52.62\% higher inflation than the tuned ordinary-benefit robust reference; the reference-minus-SAV interval favors the reference.
At Synthetic cap 2,500, tuned SAV-Val has 46.79\% higher inflation than the tuned ordinary-benefit robust reference; the reference-minus-SAV interval favors the reference.
At Synthetic cap 5,000, tuned SAV-Val has 30.09\% higher inflation than the tuned ordinary-benefit robust reference; the reference-minus-SAV interval favors the reference.

Both roles and all preset capacities remain in the associated statistics.
All nine scheduled source/cap cohorts and all 324 reference cells are independently audited, with a matching full 972-cell tuning-parent audit. The conditional,
unadjusted intervals concern the fitted models and historical validation
cohorts; they are not guarantees for new markets or seller behavior.

\subsection{Frozen-model behavioral and missingness sensitivities}
The saved-model follow-up is complete and independently audited: 225 cells
cover all nine source/seed pairs, six response or missingness changes,
four main methods, and one generic disclosure-monotone control per pair.
The same estimators and targets are retained. Concealment costs are in
stipulated log-quote utility units per hidden field. Models were fitted
with at most two hidden fields; the one- and three-hidden evaluations
change allowable reports without retraining. Added natural missingness
uses common listing/field draws across methods, and the newly unavailable
values cannot be recovered. These are retrospective sensitivities,
not new source observations or retrained robustness guarantees.

\begin{table*}[t]
\centering\small
\setlength{\tabcolsep}{4pt}
\caption{All frozen-model sensitivity settings at 10\% capacity. Cells
are positive log-quote inflation averaged over three fitting seeds.
SAV denotes full SAV-Val; robust denotes the no-feedback control.
The primary setting remains visible. Costs and added natural missingness
change stipulated behavior and available information, respectively.}
\label{tab:frozen_sensitivities}
\begin{tabular}{lrrrrrr}
\toprule
& \multicolumn{2}{c}{Synthetic} & \multicolumn{2}{c}{JUCars} & \multicolumn{2}{c}{MUCars}\\
Setting & SAV & Robust & SAV & Robust & SAV & Robust\\
\midrule
Primary: cost 0, H=2 & 0.02075 & 0.01339 & 0.01549 & 0.01154 & 0.00689 & 0.00774\\
Cost 0.02 / field & 0.00256 & 0.00494 & 0.01100 & 0.00903 & 0.00270 & 0.00281\\
Cost 0.05 / field & 0.00077 & 0.00188 & 0.00756 & 0.00583 & 0.00062 & 0.00056\\
H=1 & 0.01131 & 0.00886 & 0.01391 & 0.01015 & 0.00543 & 0.00631\\
H=3, without retraining & 0.02074 & 0.01358 & 0.01563 & 0.01154 & 0.00712 & 0.00797\\
Added natural missingness 10\% & 0.02152 & 0.01252 & 0.01447 & 0.01076 & 0.00724 & 0.00757\\
Added natural missingness 30\% & 0.01636 & 0.01054 & 0.01259 & 0.00884 & 0.00730 & 0.00684\\
\bottomrule
\end{tabular}
\end{table*}

The synthetic comparison reverses with positive concealment cost:
SAV-Val has 48.25\% and 58.85\% lower inflation than robust/no-feedback
at costs 0.02 and 0.05, with paired conditional intervals above zero.
It remains worse under both hidden-field changes and both added-missingness
settings. On JUCars, every follow-up retains higher SAV-Val inflation.
On MUCars, the one- and three-hidden settings retain positive paired
differences. Added 10\% natural missingness reduces the improvement to
4.36\%, with an interval crossing zero; at 30\%, the point comparison
reverses and its interval also crosses zero. Both positive-cost MUCars
intervals cross zero. These outcomes limit extrapolation of the primary
MUCars gain across stipulated response and missingness conditions.

\subsection{Zero uplift and valuation error}
The generic no-query disclosure-monotone closure evaluates the common
mask-aware backbone over all further withholding patterns and takes the
largest quotation. Additional withholding cannot raise that maximum.
This is a generic monotonicity diagnostic, not a novel acquisition
method or an official strategic-classification reproduction.

\begin{table*}[t]
\centering
\caption{Generic disclosure-monotone closure on 1,000 records per source,
averaged over three fitting seeds. All queries are zero. Honest log-quote
lift is relative to the common mask-aware backbone; closure bias is
relative to the asking-price target. Native-price MAEs are source-specific.
Zero uplift alone does not establish accurate or unbiased valuation.}
\label{tab:closure_tradeoffs}
\begin{tabular}{lrrrrr}
\toprule
Source & Backbone MAE & Closure MAE & Honest log lift & Closure bias & Inflation\\
\midrule
Synthetic & 3503.62 & 5346.70 & 0.09237 & 0.10010 & 0\\
JUCars & 6875.38 & 6881.40 & 0.02021 & 0.05017 & 0\\
MUCars & 27281.98 & 28386.05 & 0.04875 & 0.07466 & 0\\
\bottomrule
\end{tabular}
\end{table*}

Closure inflation is zero in every saved cell. Honest MAE increases by
1,843.08 synthetic price units, with paired conditional interval
$[1582.04,2137.71]$, and by 1,104.07 MUCars units, with interval
$[540.28,1674.32]$. The JUCars change is only 6.02 units and its interval
$[-28.85,39.34]$ crosses zero. Honest quotations rise in all sources.
The strategic metric must be assessed jointly with error and signed bias.

\subsection{Original source-missingness strata}
Technical strata count originally unrecorded optional fields before
deliberate withholding. Table~\ref{tab:technical_strata} retains every
observed group for the full and robust methods at the primary capacity.
Counts concern unique listings, not independent copies for each seed.
These groups do not constitute a demographic fairness assessment.
All four main methods, signed biases and other stratum metrics remain
in the bound source-missingness artifact.

\begin{table*}[t]
\centering\small
\setlength{\tabcolsep}{4pt}
\caption{Original source-missingness strata at 10\% capacity, with means
across three fitting seeds. Query and failed-query percentages use the
whole stratum as denominator; failure is not conditional on a query.
Missing counts concern the four optional fields, not seller actions.
Small groups remain visible and do not support population safety claims.}
\label{tab:technical_strata}
\begin{tabular}{llrlrrrrr}
\toprule
Source & Missing & $n$ & Method & Honest MAE & Attack MAE & Inflation & Query (\%) & Failed (\%)\\
\midrule
Synthetic & 0 & 863 & Robust & 3490.99 & 3583.10 & 0.01420 & 8.30 & 0.00\\
 &  &  & SAV-Val & 3475.62 & 3748.79 & 0.02245 & 8.30 & 0.00\\
Synthetic & 1 & 128 & Robust & 3806.52 & 3736.70 & 0.00820 & 21.61 & 19.53\\
 &  &  & SAV-Val & 3603.00 & 3600.97 & 0.01009 & 21.88 & 17.71\\
Synthetic & 2+ & 9 & Robust & 7584.56 & 7143.33 & 0.01034 & 7.41 & 7.41\\
 &  &  & SAV-Val & 6868.46 & 6433.93 & 0.00965 & 3.70 & 3.70\\
\midrule
JUCars & 0 & 800 & Robust & 8105.59 & 8098.47 & 0.01228 & 4.96 & 0.00\\
 &  &  & SAV-Val & 8111.44 & 8095.18 & 0.01716 & 7.08 & 0.00\\
JUCars & 1 & 200 & Robust & 2986.10 & 3000.16 & 0.00857 & 30.17 & 30.17\\
 &  &  & SAV-Val & 2978.07 & 2988.97 & 0.00882 & 21.67 & 21.50\\
\midrule
MUCars & 0 & 972 & Robust & 27470.33 & 27680.75 & 0.00772 & 9.74 & 0.00\\
 &  &  & SAV-Val & 27050.39 & 27208.00 & 0.00689 & 9.57 & 0.00\\
MUCars & 1 & 28 & Robust & 31317.13 & 31720.57 & 0.00811 & 19.05 & 9.52\\
 &  &  & SAV-Val & 30515.60 & 30920.39 & 0.00679 & 25.00 & 13.10\\
\bottomrule
\end{tabular}
\end{table*}

The two-or-more-missing synthetic group contains nine listings, and the
one-missing MUCars group only 28. These sizes limit general inference.
Failed replay concentrates in source-missing groups: for JUCars SAV-Val,
the one-missing group receives 21.67\% query attempts and 21.50\%
failed attempts. Its complete-field group has no failed replay.
This illustrates the capacity consequence of publicly indistinguishable
absence, without identifying individual sellers as strategic.

\subsection{Whole-field categorical SMI adaptation}
The full-field no-query SMI adaptation has completed all three sources.
One deterministic fitting generation per source uses all allowed covariates,
four whole optional-field groups, and the full sixteen-model linear-SVR
family. All fits share complete-optional-field estimator-training rows.
Training-only scaling, category vocabularies and target normalization are
retained. Missing categories require a one-hot completion, including the
declared unknown token. These discrete groups and fallback conventions
are a task adaptation of the two-stage SMI formulation~\cite{zhang2020strategic},
not an official Weka reproduction.

Independent categorical enumeration and scalar interval optimization check
94,527 report/domain cases, every final full-model quote, legal seller report
and source cohort. The interval oracle agrees with a separate continuous-LP
oracle on the complete execution smoke and randomized sign-consistent families.
It preserves all domains and constraints; once-only array decoding changes
audit overhead, not evaluation.

\begin{table*}[t]
\centering\small\setlength{\tabcolsep}{3pt}
\caption{Whole-field SMI adaptation, without queries, on 1,000 evaluation
listings per source. There is one deterministic fit per source, not three
independent seeds. All methods share the stated complete-optional-field
training rows. Fallback uses ordinary median/mode SVR when the fully reduced
sign-consistent family or optimization is unavailable. The original method's
guarantee is not asserted on those rows. MAEs use native source price units.}
\label{tab:grouped_smi}
\begin{tabular}{llrrrrrr}
\toprule
Source & Domain/method & Fit rows & Honest MAE & Attack MAE & Inflation & \shortstack{Honest\\fallback (\%)} & \shortstack{Attack\\fallback (\%)}\\
\midrule
Synthetic & Median/mode SVR & 2947 & 6973.99 & 7825.89 & 0.08741 & 0.0 & 0.0\\
 & SMI / free & 2947 & 6973.99 & 7825.89 & 0.08741 & 13.7 & 75.3\\
 & SMI / nonnegative & 2947 & 6973.99 & 7825.89 & 0.08741 & 13.7 & 75.3\\
 & SMI / unit box & 2947 & 6973.99 & 7825.89 & 0.08741 & 13.7 & 75.3\\
\midrule
JUCars & Median/mode SVR & 2674 & 9138.21 & 8723.06 & 0.12962 & 0.0 & 0.0\\
 & SMI / free & 2674 & 9138.21 & 8723.06 & 0.12962 & 20.0 & 87.3\\
 & SMI / nonnegative & 2674 & 9138.21 & 8723.06 & 0.12962 & 20.0 & 87.3\\
 & SMI / unit box & 2674 & 9138.21 & 8723.06 & 0.12962 & 20.0 & 87.3\\
\midrule
MUCars & Median/mode SVR & 3368 & 54377.45 & 58214.93 & 0.10496 & 0.0 & 0.0\\
 & SMI / free & 3368 & 54377.45 & 58214.93 & 0.10496 & 2.8 & 65.6\\
 & SMI / nonnegative & 3368 & 54377.45 & 58214.93 & 0.10496 & 2.8 & 65.6\\
 & SMI / unit box & 3368 & 54377.45 & 58214.93 & 0.10496 & 2.8 & 65.6\\
\bottomrule
\end{tabular}
\end{table*}

All three adapted domains have the same aggregate honest and attacked
outcomes as ordinary median/mode SVR. Attacked fallback fractions are
75.3\%, 87.3\% and 65.6\% on synthetic, JUCars and MUCars respectively.
This demonstrates the importance of retained-family assumptions when
extending the method to categorical market representations. It does not
demonstrate failure of the original theorem under its assumptions, or a
gain attributable specifically to SAV-Val acquisition. The numerical and
categorical diagnostics have different fitting rows and action spaces.

\subsection{Acquisition-conditioned neighbor control}
ACO-1 is a clean-room regression/group/hard-cap adaptation of
acquisition-conditioned expected-loss scoring~\cite{valancius2024aco}.
Each price predictor is fixed. A training-only conditional neighbor bank
provides squared log-loss changes for sixteen optional-field observation
patterns; unavailable bank values give zero replay benefit. Distances use
only observed query groups and all allowed nonoptional groups, with fixed
numeric/category weighting and stable ties. One-field acquisition and
batch capacity replace general multi-step rollout and the acquisition-cost
trade-off. This is not an unchanged execution of the authors' classification system.

The complete nine-pair study contains 144 cells. Main $k=25$ uses all six
capacities; $k=10$ and $k=50$ use only the primary capacity. Independent
checks cover role separation, loss banks, field boundaries, direct coordinate
distances, allocations, sixteen metrics and first-batch exact unilateral
deviations. No winning $k$ is selected.

\begin{table*}[t]
\centering
\caption{All ACO-1 neighbor settings at 10\% capacity, averaged over three
fitting seeds. Mask and robust are the common mask-aware and robust/no-feedback
price predictors. Native-price MAEs are source-specific. Main $k=25$ and
both fixed sensitivity counts remain visible.}
\label{tab:aco_primary}
\begin{tabular}{llrrrr}
\toprule
Source & Price predictor & $k$ & Inflation & Honest MAE & Attack MAE\\
\midrule
Synthetic & Mask & 10 & 0.06439 & 3503.62 & 4624.94\\
 &  & 25 & 0.05742 & 3503.62 & 4416.39\\
 &  & 50 & 0.05684 & 3503.62 & 4412.09\\
 & Robust & 10 & 0.01089 & 3568.22 & 3618.77\\
 &  & 25 & 0.01081 & 3568.22 & 3612.31\\
 &  & 50 & 0.01042 & 3568.22 & 3609.03\\
\midrule
JUCars & Mask & 10 & 0.01762 & 6875.38 & 6880.10\\
 &  & 25 & 0.01779 & 6875.38 & 6881.85\\
 &  & 50 & 0.01819 & 6875.38 & 6881.19\\
 & Robust & 10 & 0.01128 & 7081.70 & 7082.40\\
 &  & 25 & 0.01118 & 7081.70 & 7081.03\\
 &  & 50 & 0.01088 & 7081.70 & 7081.48\\
\midrule
MUCars & Mask & 10 & 0.03838 & 27281.98 & 28064.34\\
 &  & 25 & 0.03817 & 27281.98 & 28083.47\\
 &  & 50 & 0.03711 & 27281.98 & 27971.09\\
 & Robust & 10 & 0.00715 & 27578.04 & 27776.84\\
 &  & 25 & 0.00716 & 27578.04 & 27778.54\\
 &  & 50 & 0.00708 & 27578.04 & 27786.83\\
\bottomrule
\end{tabular}
\end{table*}

With robust prices and $k=25$, ACO-1 has lower primary inflation than full
SAV-Val on synthetic and JUCars records. Full-minus-ACO differences are
$0.00995$ and $0.00431$, with conditional intervals
$[0.00725,0.01345]$ and $[0.00291,0.00588]$. On MUCars, the point difference
favors SAV-Val by $0.00027$, but its interval crosses zero. ACO-1 also
improves on robust suppression-output acquisition on synthetic and MUCars
records, while the JUCars interval crosses zero. These results strengthen
the need for a strong acquisition control; the primary comparison varies by source.

\subsection{Shared-report transfer}
The original main run also saved outcomes when every principal method
receives full SAV-Val's returned primary reports. Each allocation and quote
is recomputed for its own fixed model and rule while reports stay fixed.
This is one-direction transfer; sellers are not reoptimized against receiving
methods. Independent reconstruction checks public head inputs, sorted
allocation and twelve reported metrics for all 36 source/seed/method cells.

\begin{table*}[t]
\centering
\caption{Principal methods on full SAV-Val's primary reports, at 10\%
capacity, averaged over three fitted seeds. The full method's own outcome
is unchanged; the others are transfer evaluations, not equilibrium outcomes
for the receiving methods. Native-price attacked MAE remains visible.}
\label{tab:shared_reports}
\begin{tabular}{llrr}
\toprule
Source & Receiving method & Inflation & Attack MAE\\
\midrule
Synthetic & Mask + benefit & 0.04355 & 4415.82\\
 & Robust, no feedback & 0.01003 & 3696.09\\
 & SAV + benefit output & 0.01768 & 3772.35\\
 & SAV-Val & 0.02075 & 3754.03\\
\midrule
JUCars & Mask + benefit & 0.00593 & 6877.03\\
 & Robust, no feedback & 0.00719 & 7110.27\\
 & SAV + benefit output & 0.01201 & 7061.26\\
 & SAV-Val & 0.01549 & 7073.94\\
\midrule
MUCars & Mask + benefit & 0.01393 & 27548.09\\
 & Robust, no feedback & 0.00473 & 27751.89\\
 & SAV + benefit output & 0.00517 & 27067.27\\
 & SAV-Val & 0.00689 & 27311.95\\
\bottomrule
\end{tabular}
\end{table*}

Robust/no-feedback has lower inflation than full SAV-Val on these shared
reports in all three sources, with reference-minus-full conditional intervals
below zero. The MUCars advantage in the own-response main study therefore
does not persist in this transfer. MUCars attacked MAE still favors SAV-Val:
the robust-minus-full difference is 439.95 price units, with interval
$[146.51,778.31]$. This separates comparison on fixed reports from comparison
under different induced responses; it does not identify a causal feedback effect.

\subsection{Online cost and stored model size}
The public inference path encodes and quotes the current report, scores
publicly absent fields, allocates queries, replays recorded values, and
revalues the batch. Model/data loading and seller search are excluded.
Private counterfactual enumeration is used by the independent auditor,
not the timed pipeline. All 360 batches reproduce parent queries, replay
outcomes and final quotes. There are 7,200 measurements, using two warmups
and twenty repetitions per batch with four CPU threads.

\begin{table*}[t]
\centering
\caption{Public inference on a shared workstation: 100 listings per batch,
three saved fitting seeds and ten batches per seed. Each timing cell has
600 repeated measurements, not independent evaluation samples. Times
include preprocessing, scoring, allocation, software replay and final
revaluation. Stored size averages compressed model-and-head files over
three seeds. These are not physical inspection times or controlled speedups.}
\label{tab:online_resources}
\begin{tabular}{llrrr}
\toprule
Source & Method & Median batch (ms) & Batch p90 (ms) & Stored size (MiB)\\
\midrule
Synthetic & Mask + benefit & 55.75 & 81.00 & 15.44\\
 & Robust, no feedback & 55.38 & 69.05 & 16.92\\
 & SAV + benefit output & 55.50 & 72.91 & 17.69\\
 & SAV-Val & 54.86 & 75.95 & 16.95\\
\midrule
JUCars & Mask + benefit & 64.19 & 78.51 & 1.98\\
 & Robust, no feedback & 65.24 & 100.76 & 1.94\\
 & SAV + benefit output & 65.17 & 77.43 & 2.01\\
 & SAV-Val & 60.33 & 68.97 & 1.99\\
\midrule
MUCars & Mask + benefit & 56.25 & 60.62 & 0.87\\
 & Robust, no feedback & 57.35 & 63.48 & 1.15\\
 & SAV + benefit output & 59.60 & 65.70 & 1.14\\
 & SAV-Val & 57.74 & 63.00 & 1.15\\
\bottomrule
\end{tabular}
\end{table*}

Hardware contention and sequential method order limit speed comparisons.
Separately recorded outer/deployment fit stages cover the shared twelve-method
matrix. Summed stage times are 316.33, 1,669.22 and 5,671.34 seconds for
synthetic, JUCars and MUCars. These do not attribute training cost to one
method, measure peak memory or energy, or include all head reconstruction,
response evaluation and artifact I/O. Nested fitting and equal tuning add
offline work that must be weighed against source-dependent gains.


\section{Common-profile deviation generalization}
\label{sec:cdm_deviation_generalization}
This exploratory diagnostic was specified after the original CDM outcomes
were available and before its own tail summaries were opened. It neither
refits the price estimator nor changes the primary positive-inflation
endpoint. All 108 source/seed/price/capacity cases and 1,512 method cells
are retained. For each held-out batch, every method faces the same five
peer-report profiles: truthful, no-query maximizing, fixed-seed random,
suppression-response and benefit-response. These constructions match the
calibration profile types but use held-out listings; they do not cover
all possible peer configurations.

For listing $i$ and peer profile $p$, we evaluate every legal report $h$,
reallocate the complete batch under both $h$ and the truthful report, and
compute the signed log-quote difference under each rule $k$, denoted
$D_{i,p,h,k}$. Given rule probabilities $\alpha$, the envelope is
\begin{equation}
 e_{i,p}(\alpha)=\max\{0,\max_h\sum_k\alpha_kD_{i,p,h,k}\}.
\label{eq:common_profile_envelope}
\end{equation}
The rule expectation in Eq.~\eqref{eq:common_profile_envelope} precedes
report maximization. Listings and profiles
receive equal weight, irrespective of how many reports are legal. We
retain the mean, maximum, 95th percentile and exact upper 5\% probability
mass mean (conditional value at risk, CVaR), using fractional boundary mass
where necessary. Whole-case
tails have descriptive seed means and standard deviations. The intervals
in Table~\ref{tab:cdm_deviation_generalization} apply instead to batch-level
measures: paired seed contrasts are averaged before resampling the ten
shared batches. They are conditional on the three fitted models and are
unadjusted exploratory intervals. ``Chosen pure'' is selected on calibration
rows, never on these held-out summaries. The complete exported contrasts
retain all four references and all six capacities.

With robust pricing at 10\% capacity, CDM reduces batch CVaR relative to
the calibration-selected pure rule by 7.34\% on Synthetic, 2.86\% on JUCars
and 6.81\% on MUCars, with positive conditional intervals. Its batch maxima
also have positive intervals against this reference on Synthetic and MUCars;
the JUCars interval includes zero. These gains do not establish an advantage
over risk: its CVaR intervals against risk include zero on all three robust
cohorts, and mean envelopes are higher on Synthetic and MUCars. With
mask-aware pricing on MUCars, mean, 95th-percentile and CVaR envelopes are
higher than risk, with intervals excluding zero. Thus finite-profile mixing
has a bounded effect against its training-selected pure reference, while
the original absence of a consistent risk-rule advantage remains.

\begin{table*}[t]
\centering
\caption{Complete common-profile deviation diagnostic at 10\% capacity.
Columns report relative reductions in batch measures (\%); positive values
favor CDM. The final column is reference minus CDM batch CVaR, in log-quote
units, with its conditional 95\% interval. Chosen pure is selected on
calibration data. Every source, estimator and reference is retained.}
\label{tab:cdm_deviation_generalization}
\setlength{\tabcolsep}{3pt}
\begin{tabular}{lllrrrrl}
\toprule
Source & Pricing & Reference & Mean & Maximum & P95 & CVaR & CVaR difference [interval]\\
\midrule
Synthetic & Masking & Risk & 1.29 & 1.87 & -0.34 & -0.33 & $-0.00087\;[-0.00330,0.00157]$\\
Synthetic & Masking & Benefit & 0.02 & 0.05 & 0.07 & 0.06 & $0.00017\;[0.00013,0.00022]$\\
Synthetic & Masking & Uniform & 11.29 & 2.13 & 1.18 & -1.03 & $-0.00266\;[-0.00876,0.00424]$\\
Synthetic & Masking & Chosen pure & 1.35 & 4.15 & 1.21 & 1.47 & $0.00391\;[0.00109,0.00726]$\\
Synthetic & Robust & Risk & -1.20 & 2.72 & 0.80 & 1.75 & $0.00143\;[-0.00010,0.00292]$\\
Synthetic & Robust & Benefit & -1.46 & 2.72 & 0.40 & 2.51 & $0.00207\;[0.00071,0.00331]$\\
Synthetic & Robust & Uniform & 9.48 & 3.89 & 19.12 & 7.28 & $0.00632\;[0.00344,0.00854]$\\
Synthetic & Robust & Chosen pure & 3.40 & 5.05 & 10.17 & 7.34 & $0.00637\;[0.00470,0.00786]$\\
JUCars & Masking & Risk & -0.56 & -0.78 & 0.63 & -1.12 & $-0.00135\;[-0.00391,0.00061]$\\
JUCars & Masking & Benefit & 4.99 & 3.97 & 2.39 & 6.18 & $0.00807\;[0.00208,0.01528]$\\
JUCars & Masking & Uniform & -1.19 & 2.67 & -3.26 & 0.62 & $0.00077\;[-0.00533,0.00801]$\\
JUCars & Masking & Chosen pure & 0.21 & -0.90 & 1.31 & 0.58 & $0.00071\;[0.00006,0.00145]$\\
JUCars & Robust & Risk & 3.34 & 3.44 & 3.03 & 3.23 & $0.00357\;[-0.00049,0.00707]$\\
JUCars & Robust & Benefit & 0.80 & -0.44 & 3.52 & -1.31 & $-0.00138\;[-0.00778,0.00494]$\\
JUCars & Robust & Uniform & -3.62 & -4.33 & -4.99 & -8.60 & $-0.00848\;[-0.01149,-0.00501]$\\
JUCars & Robust & Chosen pure & 4.77 & 2.46 & 3.04 & 2.86 & $0.00315\;[0.00141,0.00524]$\\
MUCars & Masking & Risk & -15.24 & -2.53 & -4.85 & -4.21 & $-0.00831\;[-0.01115,-0.00586]$\\
MUCars & Masking & Benefit & 4.43 & 7.66 & 1.84 & 3.59 & $0.00765\;[0.00071,0.01520]$\\
MUCars & Masking & Uniform & -1.63 & 4.01 & -3.74 & -0.78 & $-0.00158\;[-0.00821,0.00574]$\\
MUCars & Masking & Chosen pure & 3.81 & 5.49 & 2.48 & 3.01 & $0.00638\;[0.00285,0.01061]$\\
MUCars & Robust & Risk & -3.67 & 1.77 & -2.82 & -1.62 & $-0.00101\;[-0.00279,0.00084]$\\
MUCars & Robust & Benefit & 8.52 & 6.07 & 6.85 & 6.39 & $0.00434\;[0.00292,0.00602]$\\
MUCars & Robust & Uniform & -5.34 & -1.63 & -8.68 & -9.00 & $-0.00525\;[-0.00616,-0.00441]$\\
MUCars & Robust & Chosen pure & 6.75 & 3.70 & 8.42 & 6.81 & $0.00465\;[0.00351,0.00601]$\\
\bottomrule
\end{tabular}
\end{table*}


\section{Discussion and Limitations}
\subsection{What the current comparison establishes}
The completed evidence demonstrates that a method can improve on a basic
masked predictor and still lose to a simpler robust control in two sources,
while improving on that robust control in a third. It does not establish
why the feedback loop's effect differs by source. Approximate training
responses, target distribution changes, and fixed hyperparameters are
possible explanations to investigate, not experimentally confirmed
causes. The current full design should therefore remain a candidate
pending prospective source evaluation and final evidence review; completed
source/size retraining does not establish a general gain.

Positive quote inflation alone also has a structural limitation. For
any reported record, taking the maximum raw quote over all further
withholding patterns produces a disclosure-monotone quote: additional
withholding shrinks that set, so it cannot raise the maximum. This can
eliminate measured positive inflation while raising honest quotes and
honest error. The completed generic envelope control exposes that
trade-off in Table~\ref{tab:closure_tradeoffs}. It is not claimed as a
novel algorithm or a prior strategic-classification implementation.

\subsection{Limits of source and behavioral evidence}
The public outcomes are asking prices. A model that predicts them
accurately may reproduce sellers' pricing preferences rather than
estimate realized sale value. Source records do not provide observed
concealment labels, query outcomes from physical inspections, or
measured seller utilities. Restoring a recorded field is an artificial
replay; original values may themselves be inaccurate. The stipulated
attacker cannot fabricate values, change the listing identifier, choose
its peers or manipulate other recorded features. Results do not extend
automatically to those behaviors or to other second-hand product classes.

The processed public validation cohorts were used during historical
development. Their retrospective status remains even if current fits
exclude evaluation labels. Capped samples and ten batches per source
also restrict representativeness and uncertainty estimation. The
training/evaluation duplicate control lacks original advertisement
identifiers. These limitations call for prospective independent data
and real query costs before marketplace deployment.

\subsection{Comparator scope and completed retraining}
Equation-level and full-field SMI adaptations are independently audited,
with explicit domains, discrete groups and fallback boundaries. The full
ACO-1 adaptation, saved-model sensitivities, technical strata, common-report
transfer and public inference diagnostic are also independently checked.
Ordinary acquisition controls remain separate from these declared
adaptations; none is an official unchanged classification reproduction.
Equal-allowance tuning and size retraining cover all 972 evaluation cells.
The completed fixed-setting ordinary-benefit robust reference and its separate
allowance extension are reported explicitly. Benefit-specific selection and
refitting cover all 324 cells, with all 27 cases and 54 deployment heads
independently checked. The completed comparisons remain heterogeneous.
Prospective source evaluation, real query costs and observed seller utility
are outside the available evidence. Nested target generation adds offline
cost that must be weighed against any measured acquisition gain.

\section{Provisional Conclusion}
We implement and audit a budgeted verification protocol with publicly
observable query eligibility, capacity-consuming failed replay, nested
target generation and seller responses to the deployed allocation.
The complete main study favors a simpler no-feedback control on strategic
quote inflation in two sources and favors SAV-Val on MUCars;
attacked-error differences also depend on the source. The full algorithm's
general empirical increment remains unresolved. The follow-up evidence
does not support a universal advantage, deployment
recommendation or journal-acceptance claim.
